\documentclass[11pt]{article}
\usepackage{amsmath,amssymb,amsthm,mathtools,hyperref}

\begin{document}
\section*{Multiscale closure}
\label{prob:042}
\label{rank:23}
\par\smallskip

\noindent\textbf{Domain.} Coarse-graining biological dynamics. \quad \textbf{Retained status/source context.} Open application-specific closure and error-control programme. Exact projected dynamics can contain memory; the supplied whole-cell review supports multiscale challenges rather than proving a universally unknown Markovian closure theorem.
\paragraph{Definitions and assumptions.} Let $x_t$ be a specified microscopic Markov process with generator $\mathcal L$ and $X_t=\Pi(x_t)\in\mathbb{R}^d$ measured coarse variables, such as metabolite amounts and cell volume. Select a finite family of observables $a_j(X)$, initial ensembles $\mathcal P_0$ and interventions $\mathcal U$. Candidate closures are Markov diffusions or finite-memory equations with parameterized drift $F$, kernel $K$ and noise. Microscopic simulations used for calibration have declared error and are separated from experimental validation. Coarse variables and memory horizon are choices to justify, not arbitrary functions to recover.
\paragraph{Problem statement.} Find a computationally cheaper closure, with approximate coarse process $\widehat X$, preserving declared observables. The displayed colored-forcing ansatz governs $\widehat X$, while $X_t=\Pi(x_t)$ remains the reference projection:
\[
\dot{\widehat X}(t)=F_\theta(\widehat X(t))+\int_0^tK_\theta(t-r,\widehat X(r))\,\,\mathrm{d} r+\eta(t),
\quad\sup_{P_0,u,0\leq t\leq T,j}\frac{|\mathbb{E} a_j(X_t)-\mathbb{E} a_j(\widehat X_t)|}{\varepsilon_j}\leq1.
\]
State conditions on microscopic mixing and projection under which the bound holds. In words, retain the information needed for reliable cell-scale prediction without simulating every molecule.
\paragraph{Additional binding assumptions and resolution criterion.} Fix the initial conditional microscopic ensemble associated with each coarse state and the law, covariance and initial-state dependence of $\eta$. The displayed ordinary derivative requires an integrable colored forcing; white noise instead requires a stochastic integral formulation. A time-convolution kernel is an ansatz, not a general exact closure for every nonlinear projection. The corrected bound normalizes each observable's tolerance separately.
For each concrete instance, specify how every $u\in\mathcal U$ acts on the microscopic dynamics and the closure, including its forcing law; whether interventions are prescribed inputs or feedback policies; and the information available to any feedback policy. Specify the coarse initialization for every $P_0\in\mathcal P_0$, its relation to the projected initial law, and whether all admitted microscopic initial laws must share the fixed conditional ensemble. These initialization and intervention conventions are not selected here. The forcing specification includes its joint temporal law and dependence on initialization. The integral from $0$ to $t$ alone does not specify a bounded memory horizon: a finite-memory instance must declare its horizon, kernel support or integration window, any required history initialization, and how truncation error enters the displayed tolerance. The intended memory convention remains to be specified. Define ``computationally cheaper'' before evaluating candidates by declaring the microscopic baseline, common prediction task and numerical accuracy, resource measure and comparison rule. Account for initialization, kernel evaluation and noise generation, and state whether calibration and preprocessing costs are included or amortized and over which workload. No particular cost measure or accounting convention is selected here.

\paragraph{Alternative formulations.} (i) Assume fast microscopic mixing conditional on $X$ and derive a memory-free closure with an explicit timescale-separation error. (ii) When slow hidden states persist, augment $X$ or retain a finite-memory kernel and test lagged correlations. A projected Markov process need not itself be Markov; the alternatives are equivalent only under additional closure conditions.
\paragraph{Obstructions and partial progress.} Multiscale models already combine molecular and cellular representations. Their agreement on means does not guarantee correct response or noise correlations. Intervention changes can destroy a separation calibrated at baseline; unobserved conserved quantities may make a chosen projection inadequate. A successful answer can establish impossibility for a specified coarse-variable set and identify the additional state needed. This differs from assembling all modules in \href{https://we-solve.github.io/problems/predictive-whole-cell-modelling-fd60530f/}{Problem~41}.
\paragraph{Sources (supplied provenance; not all rechecked).}
\begin{itemize}
\item Georgouli et al., Multi-scale models of whole cells: progress and challenges (2023; review; publisher/indexed versions checked after access challenge). \url{https://pmc.ncbi.nlm.nih.gov/articles/PMC10661945/}
\item Georgouli et al., Multi-scale models of whole cells: progress and challenges (2023; accessible publisher version; biological motivation). \url{https://www.frontiersin.org/journals/cell-and-developmental-biology/articles/10.3389/fcell.2023.1260507/full}
\end{itemize}


\section*{Binding catalogue conventions}
This edition revises both supplied catalogues and preserves original problem numbers 1--330. It adds one globally unique label \verb|\label{rank:x}| per problem, where smaller $x$ means greater provisional difficulty. A rank label is an anchor at that problem's section: \verb|\ref{rank:x}| returns the section number, not the literal difficulty number. The visible \texttt{rank:x} text and the index give the difficulty number itself.

The entries are \emph{parameterized mathematical research schemes}, including empirical prediction, causal identification, historical reconstruction, theory and engineering. They are not 330 certified unsolved mathematical conjectures. A rigorous instance fixes the objects and restrictions below; changing those objects changes the question. Established identities, finite calculations, non-identification constructions and solved special cases are not being posed anew as open. Some residual tasks concern informative bounds, biological realization or external validation rather than a new theorem.

Every main definition and problem statement received a mathematical editorial review. Each entry now has an additional binding-conditions paragraph; direct corrections are separately recorded in the accompanying audit. The alternative formulations inherit the applicable conventions, but are different targets unless equivalence is proved. Where an old broad phrase conflicts with an explicit restriction below it, the explicit restriction controls.

\textbf{Literature boundary.} Thirteen selective primary-literature checks cover eighteen entries in this revision; their access levels and scope are recorded explicitly. The supplied bibliography and its earlier source-access descriptions are retained as provenance, not represented as newly verified. Neither this review nor a source's description proves the absence of later solutions. The date of this edition is 13 October 2026; there is no claim of exhaustive current-status certification for all 330 entries.

\subsection*{Binding conventions for a concrete instance}
\paragraph{Objects, quantifiers and non-vacuity.}
An instance declares measurable state, input, output and parameter spaces; a fixed candidate family; admissible initial laws, interventions and environments; a finite horizon or specified asymptotic regime; observation and sampling laws; and performance metrics, positive normalization scales and positive tolerances. These are chosen before fitting or testing, not adjusted after seeing failures. Shared parameters must genuinely be shared across the declared validation domain. Finite dimension alone does not restrict arbitrary real-number encodings: specify a finite library, bounded description length, or regularity/complexity bounds. Exclude constant or no-signal solutions through a fixed informative-output or baseline-improvement requirement where appropriate. A universally quantified result cannot be established by choosing one trivial favorable instance.

\paragraph{Probability laws and data.}
Every expectation or probability refers to the specified stochastic law, including initial uncertainty, process noise, observation noise and sampling as appropriate. A deterministic equation is not a probability model without a declared random input or initial law. ``Observed'' targets in population-level formulas mean population observables $g_j^*$, not noisy estimates asserted to be exact. Let $\widetilde g_j$ be independent validation estimates and suppose
\[
\mathbb{P}\{\lvert\widetilde g_j-g_j^*\rvert\le r_j\text{ for every }j\}\ge1-\alpha.
\]
Then an uncertainty-aware sufficient certification rule for positive target tolerances $\varepsilon_j$ is
\[
\max_j\frac{|\widehat g_j-\widetilde g_j|+r_j}{\varepsilon_j}\le1.
\]
On that simultaneous confidence event it implies every population error is at most its tolerance, by the triangle inequality. Vector/norm versions use corresponding valid simultaneous bounds. Training uncertainty must also be included when the claim averages over repeated training data. Evaluation on finitely many conditions does not establish a continuous-domain supremum without covering and regularity bounds.

\paragraph{Dynamics and well-posedness.}
For ODEs and SDEs give coefficient domains, initial laws and existence/uniqueness conditions in the intended solution sense, or keep the entire admitted solution set. For stochastic jump models specify nonnegative off-diagonal rates, diagonal normalization, rate orientation and non-explosion. For PDEs specify spatial domains, boundary/initial conditions, conservation or source terms, and a weak or classical solution class. Moving-domain concentration equations include dilution/transport terms. A stationary law is used only for a time-homogeneous augmented model under existence assumptions; a stationary approximation to a driven system must be justified. Absorbing extinction states are retained. Survival-conditioned, quasi-stationary and unconditional laws are distinct and must not be substituted for one another.

\paragraph{Identification and causal meaning.}
Two parameter values are observationally equivalent when they induce the same full distribution under all allowed observation/intervention designs, not merely the same likelihood value for one sample. A target is identified only if constant on the equivalence class. Gauge symmetries, label switching and reference conventions are quotiented or fixed. Finite-data compatible sets use a declared confidence, likelihood or discrepancy rule; an empty set signals model rejection. Bayesian credible sets are not automatically confidence sets. A prior or regularizer does not create data identification. Bounds intended to resolve a question must be informative at a predeclared width or decision threshold.

A causal contrast uses well-defined intervention versions and a stated causal model. Random assignment, or explicit consistency, overlap and exchangeability assumptions, is required for causal interpretation from observations. Interference and time-varying confounding are represented when present. Fixing a mediator by intervention differs from conditioning on its observed value. Marginal potential-outcome laws do not determine individual cross-world effects or an arbitrary joint coupling. Ethically admissible designs, rather than unconstrained interventions, define the experimental menu.

\paragraph{Metrics, units, optimization and limits.}
Specify the sampling or reference measure for densities and integrals, the log base, and dimensionless arguments of logarithms. Total variation is $\sup_A|P(A)-Q(A)|$; Wasserstein distances require the declared ground metric and appropriate moment conditions; KL divergence uses its stated absolute-continuity convention. All normed readouts are scaled consistently, and error probabilities use $0<\alpha<1$. A ratio has a positive denominator or an explicit undefined/zero convention. Thermodynamic entropy production is nonnegative in expectation under the relevant hypotheses; individual stochastic realizations need not be. An infimum is not an attained minimum without compactness/coercivity and continuity or another existence argument. Fix the order of large-population, long-time and vanishing-mutation limits. Stability claims declare perturbation space and symmetries.

\subsection*{How the provisional difficulty order was assigned}
Four editorial components are scored from 1 (lower burden) to 5 (higher burden): (a) formal or inferential difficulty, (b) access to evidence and identification burden, (c) mechanistic/computational integration, and (d) independent validation and transfer. Their unweighted total orders the entries. For selected entries in the highest bands, an explicit editorial priority breaks equal totals; all other ties are broken by original problem number. Equal totals are \emph{not} evidence of reliably different true difficulties. Scores are judgments, not measurements, completion-time estimates or consensus rankings.

The object ranked is an informative, nondegenerate resolution of the \emph{scoped} scientific target, with the stated budget and validation domain. A sharp useful impossibility result can qualify; an already known identity, vacuous bound, arbitrary favorable model or lookup table cannot. Rank 330 is merely the least difficult relative assessment in this catalogue, not a claim of an easy biological solution. Broad historical titles do not automatically outrank precisely formulated mechanistic tasks.



\end{document}
